\section{Fibre degree and the cyclic-six centre}
\label{zdc:sec:bridge}

This section supplies the theorem--proof module behind the fibre-degree and
cyclic-centre formulas recorded in the orientation material.  There are three
different objects in the calculation:
\[
 H^4(U;\mathbb Z),\qquad H^4(U;\mathbb Z/6),\qquad
 Z(SU(6))=\mu _6.
\]
The first is a free rank-six characteristic lattice, the second has six
independent residue coordinates, and the third is a single cyclic group.
The candidate determines one particular surjection from the first object to
the third.  It does not identify all three objects, and it does not construct
a gauge connection or a spectral operator.

\subsection{Spaces, markings, and the finite-band presentation}
\label{zdc:subsec:markings}

Retain the reconstructed map
\[
 f:X\longrightarrow\mathbb P^1,
 \qquad X\cong_{\mathrm{diff}}S^6,
\]
and write
\[
 D_0=W,\qquad D_1=S_1,\qquad D_2=S_2,
 \qquad D=D_0\sqcup D_1\sqcup D_2,
 \qquad U=X\setminus D.
\]
Here $W$ is the reduced non-normal fibre and $3S_1,4S_2$ are the
two multiple fibres.  The verified local homology computations give
\[
 H_1(W;\mathbb Z)=\mathbb Z\langle a_0,b_0\rangle,
 \qquad
 H_1(S_j;\mathbb Z)=\mathbb Z\langle a_j,b_j\rangle
 \quad(j=1,2).
\]
The marking is
\[
 a_0=[\widehat\gamma],\quad b_0=[\widehat u],
 \quad a_j=[\widehat g_j],
\]
where $\widehat g_j$ is the clockwise affine deck generator.  The second
generator $b_j$ is the invariant coinvariant coordinate.  With the
marking used in the finite-band calculation one may take
\(b_1=\pi_{1*}[\widehat w]\) and
\(b_2=\pi_{2*}[\widehat w]\); in the second coinvariant quotient
\(\pi_{2*}[\widehat u]=\pi_{2*}[\widehat w]\), so this agrees with the
alternative representative used elsewhere.  The two specializations of
the marked $\widehat\gamma$-cycle are
\begin{equation}
 \pi_{1*}[\widehat\gamma]=3a_1,
 \qquad
 \pi_{2*}[\widehat\gamma]=-4a_2.
 \label{zdc:eq:gamma-specializations}
\end{equation}

Fix Alexander duality using the complex orientation of $X$, outward-first
boundary orientation, and the convention that the Poincare dual of a
positive nearby fibre is the negative covector of its clockwise meridian.
It identifies
\begin{equation}
 L:=H^4(U;\mathbb Z)
 \xrightarrow{\ \sim\ }
 H_1(D;\mathbb Z)
 =\mathbb Z\langle a_0,b_0;a_1,b_1;a_2,b_2\rangle
  \cong\mathbb Z^6.
 \label{zdc:eq:defect-lattice}
\end{equation}
Indeed, $D$ is compact and locally contractible, so
\[
 \widetilde H^q(S^6\setminus D;\mathbb Z)
 \cong\widetilde H_{5-q}(D;\mathbb Z).
\]
Take $q=4$ and use the disjoint-union decomposition of $D$.

Let $X^\circ$ be the space obtained after the two finite fillings and
before the central toric filling.  For logarithmic twist integers
\(\ell_1,\ell_2\), its verified degree-one Mayer--Vietoris calculation is
the presentation
\begin{equation}
 H_1(X^\circ;\mathbb Z)
 =\frac{\mathbb Z\langle c,\bar g_1,\bar g_2\rangle}
 {\mathbb Z\langle3\bar g_1-\ell_1c,
                         4\bar g_2-\ell_2c\rangle}.
 \label{zdc:eq:h1xcirc-presentation}
\end{equation}
Thus $c,\bar g_1,\bar g_2$ are three marked generators, not a basis.
The relation matrix is
\[
 R_1=
 \begin{pmatrix}-\ell_1&3&0\\-\ell_2&0&4\end{pmatrix}.
\]
For admissible twists $3\nmid\ell_1$ and $2\nmid\ell_2$, the gcd of
its rank-two minors
\[
 3\ell_2,\qquad-4\ell_1,\qquad12
\]
is one.  Therefore the presented group is torsion-free of rank one.  The
homomorphism
\begin{equation}
 \begin{aligned}
 \Phi_1:H_1(X^\circ;\mathbb Z)&\longrightarrow\mathbb Z,\\
 \Phi_1(c)&=12,&
 \Phi_1(\bar g_1)&=4\ell_1,&
 \Phi_1(\bar g_2)&=3\ell_2
 \end{aligned}
 \label{zdc:eq:phi1}
\end{equation}
annihilates both relations.  Its three displayed values have gcd one under
the same admissibility hypotheses; hence it is an isomorphism.  For the
candidate
\begin{equation}
 (\ell_0,\ell_1,\ell_2)=(0,1,-1),
 \qquad
 \Phi_1(c)=12,\quad
 \Phi_1(\bar g_1)=4,\quad
 \Phi_1(\bar g_2)=-3.
 \label{zdc:eq:candidate-twists}
\end{equation}
Put
\begin{equation}
 \zeta:=\bar g_1+\bar g_2.
 \label{zdc:eq:zeta-generator}
\end{equation}
Then $\Phi_1(\zeta)=1$.  Since $\Phi_1$ is an isomorphism, $\zeta$ is the
positive basis element of $H_1(X^\circ;\mathbb Z)$ selected by $\Phi_1$, and
\begin{equation}
 c=12\zeta,\qquad \bar g_1=4\zeta,\qquad \bar g_2=-3\zeta.
 \label{zdc:eq:presentation-to-basis}
\end{equation}

Let
\[
 \overline q_{\mathrm{col}}:H_1(D;\mathbb Z)
 \longrightarrow H_1(X^\circ;\mathbb Z)
\]
be the homomorphism sending $a_0,a_1,a_2$ to
$c,\bar g_1,\bar g_2$, respectively, and sending every $b_i$ to zero.
Define $q_{\mathrm{col}}=\overline q_{\mathrm{col}}\circ\operatorname{AD}$,
where $\operatorname{AD}$ is the fixed Alexander isomorphism above.  Thus
the marked collar homomorphism is well typed:
\begin{equation}
 \begin{aligned}
 q_{\mathrm{col}}:L&\longrightarrow H_1(X^\circ;\mathbb Z),\\
 q_{\mathrm{col}}(a_0)&=c,&q_{\mathrm{col}}(b_0)&=0,\\
 q_{\mathrm{col}}(a_1)&=\bar g_1,&q_{\mathrm{col}}(b_1)&=0,\\
 q_{\mathrm{col}}(a_2)&=\bar g_2,&q_{\mathrm{col}}(b_2)&=0.
 \end{aligned}
 \label{zdc:eq:qcol}
\end{equation}
The central-band projection sends $a_0=[\widehat\gamma]$ to $c$ and
$b_0=[\widehat u]$ to zero.  In the two finite-band cokernels, the actual
Mayer--Vietoris map is $(\pi_{1*},-\pi_{2*})$; it sends the affine generators
to $\bar g_1,\bar g_2$ and kills the invariant coordinates $b_1,b_2$.
These statements derive the six displayed values from the marked collar
calculation.  The map is surjective because its images contain the three
presentation generators in \eqref{zdc:eq:h1xcirc-presentation}.  It is not
an inclusion-induced map $H_1(D)\to H_1(X^\circ)$: no inclusion of the
deleted central fibre $W$ into $X^\circ$ exists.

\subsection{The exact fibre-degree row}
\label{zdc:subsec:fibre-degree}

Let $F\cong T^4$ be a complex-oriented smooth fibre and choose
\(\omega_F\in H^4(F;\mathbb Z)\) with
\(\langle\omega_F,[F]\rangle=1\).  If $i_F:F\hookrightarrow U$ is the
inclusion, put $\operatorname{res}_F=i_F^*:H^4(U;\mathbb Z)\to
H^4(F;\mathbb Z)$.  Define
\begin{equation}
 \operatorname{res}_F(\alpha)=\deg_F(\alpha)\omega_F,
 \qquad
 \deg_F:L\longrightarrow\mathbb Z.
 \label{zdc:eq:degree-definition}
\end{equation}

\begin{theorem}[Fibre-degree row]
\label{zdc:thm:fibre-degree}
With the fixed Alexander-duality convention above, one has
\begin{equation}
 \boxed{
 \deg_F(n_0a_0+r_0b_0+n_1a_1+r_1b_1+n_2a_2+r_2b_2)
 =12n_0+4n_1-3n_2.}
 \label{zdc:eq:degree-row}
\end{equation}
Equivalently,
\[
 [\deg_F]_{(a_0,b_0,a_1,b_1,a_2,b_2)}
 =(12,0,4,0,-3,0).
\]
Moreover,
\begin{equation}
 \boxed{\deg_F=\Phi_1\circ q_{\mathrm{col}}}
 \label{zdc:eq:degree-factorization}
\end{equation}
for the candidate twists, and
\begin{equation}
 \boxed{
 \ker\deg_F
 =\mathbb Z\langle
 b_0,b_1,b_2,a_0-3a_1,3a_1+4a_2\rangle.}
 \label{zdc:eq:degree-kernel}
\end{equation}
This kernel is saturated of rank five.
\end{theorem}

\begin{proof}
We retain every sign in the one-shot collar calculation.  In the ordered
boundary-source basis
\[
 (\gamma_0,u_0,\widetilde a_0;
   \gamma_1,w_1,\rho_1;
   \gamma_2,w_2,\rho_2)
\]
and ordered target basis
\[
 (c,\bar\rho_1,\bar\rho_2;
   a_0,b_0;a_1,b_1;a_2,b_2),
\]
the actual map $(i_{U*},-i_{N*})$ at
\(\left(\ell_0,\ell_1,\ell_2\right)=(0,1,-1)\) is
\begin{equation}
 \Delta_1=
 \begin{pmatrix}
 1&0&0&1&0&0&1&0&0\\
 0&0&-1&0&0&1&0&0&0\\
 0&0&-1&0&0&0&0&0&1\\
 -1&0&0&0&0&0&0&0&0\\
 0&-1&0&0&0&0&0&0&0\\
 0&0&0&-3&0&-1&0&0&0\\
 0&0&0&0&-1&0&0&0&0\\
 0&0&0&0&0&0&4&0&-1\\
 0&0&0&0&0&0&0&-1&0
 \end{pmatrix}.
 \label{zdc:eq:collar-matrix}
\end{equation}
For example, the columns say
\[
 \gamma_0\mapsto c-a_0,\quad u_0\mapsto-b_0,\quad
 \widetilde a_0\mapsto-\bar\rho_1-\bar\rho_2,
\]
\[
 \gamma_1\mapsto c-3a_1,\quad w_1\mapsto-b_1,\quad
 \rho_1\mapsto\bar\rho_1-a_1,
\]
\[
 \gamma_2\mapsto c+4a_2,\quad w_2\mapsto-b_2,\quad
 \rho_2\mapsto\bar\rho_2-a_2.
\]
Thus the entries $-3,+4$ are exactly
\eqref{zdc:eq:gamma-specializations}; all other displayed minus signs are
the second-inclusion signs in the Mayer--Vietoris convention.  Direct
integral elimination gives $\det\Delta_1=1$.

Let $\widetilde a_0^\vee$ be the integral covector on the ordered
boundary-source group that is one on $\widetilde a_0$ and zero on the other
eight basis elements.  Pullback of target covectors is represented by
$\Delta_1^{\mathsf T}$.  By the fixed orientation convention, the target
covector extending the Poincare dual of a complex-oriented nearby fibre
must pull back to $-\widetilde a_0^\vee$.  Since $\Delta_1$ is unimodular,
there is a unique integral extension.  Its last six entries are precisely
the Alexander coordinates of the restriction functional on $H_1(D)$.
Consequently they are obtained from
\begin{equation}
 \Delta_1^{\mathsf T}y=-\widetilde a_0^\vee.
 \label{zdc:eq:transpose-system}
\end{equation}
Writing
\[
 y=(A,B,C;X_0,Y_0;X_1,Y_1;X_2,Y_2),
\]
the nine equations are
\[
 A=X_0,\quad Y_0=Y_1=Y_2=0,\quad
 A=3X_1,\quad A=-4X_2,
\]
\[
 B=X_1,\qquad C=X_2,\qquad -B-C=-1.
\]
Their unique solution is
\[
 y=(12,4,-3;12,0;4,0;-3,0).
\]
The last six entries prove \eqref{zdc:eq:degree-row}.  Equations
\eqref{zdc:eq:phi1}, \eqref{zdc:eq:candidate-twists}, and
\eqref{zdc:eq:qcol} give the same row on the six basis elements, proving
the typed equality \eqref{zdc:eq:degree-factorization}.

The row is primitive because $\gcd(12,4,3)=1$; explicitly,
\(\deg_F(a_1+a_2)=1\).  Each of the five displayed elements in
\eqref{zdc:eq:degree-kernel} lies in the kernel.  Conversely, if
\(12n_0+4n_1-3n_2=0\), reduction modulo four gives $n_2=4v$, and then
$n_1=3v-3n_0$.  Therefore
\[
 n_0a_0+n_1a_1+n_2a_2
 =n_0(a_0-3a_1)+v(3a_1+4a_2).
\]
The three $b_i$-coordinates are free, and the five displayed generators
are independent.  They therefore form a basis of the kernel.  Finally, a
kernel of a surjection from a free abelian group to $\mathbb Z$ is
saturated, proving the last assertion.
\end{proof}

Reversing the global Alexander-duality convention negates all six entries
of \eqref{zdc:eq:degree-row} simultaneously.  It cannot change the sign of
only one component.  In particular, with the convention fixed here,
$q_{\mathrm{col}}(a_2)=+\bar g_2$ while the coordinate of $\bar g_2$ in the
positive generator of $H_1(X^\circ;\mathbb Z)$ is $-3$.

For completeness, if $j:M\hookrightarrow X^\circ$ denotes the central
collar inclusion, the toric attaching meridian obeys
\[
 j_*[\widetilde a_0]=\ell_0c-\bar g_1-\bar g_2,
\]
and hence
\[
 \Phi_1(j_*[\widetilde a_0])
 =12\ell_0-4\ell_1-3\ell_2=p.
\]
At the candidate values this is $p=-1$.  Thus the same correctly typed
functional controls the smooth-fibre degree and the terminal meridian, but
the two statements use the maps displayed above rather than an untyped
identification of their source groups.

\subsection{Normal-line holonomy and the two finite reductions}
\label{zdc:subsec:finite-holonomy}

The transverse disc in the order-$m_j$ local quotient is acted on by
\(e^{-2\pi i/m_j}\), with $m_1=3,m_2=4$.  Therefore the normal-line
holonomy characters, in the marked bases, are
\begin{equation}
 \begin{aligned}
 h_1:H_1(S_1;\mathbb Z)&\longrightarrow\mu_3,&
 h_1(a_1)&=e^{-2\pi i/3},& h_1(b_1)&=1,\\
 h_2:H_1(S_2;\mathbb Z)&\longrightarrow\mu_4,&
 h_2(a_2)&=e^{-2\pi i/4},& h_2(b_2)&=1.
 \end{aligned}
 \label{zdc:eq:normal-holonomies}
\end{equation}
Via \eqref{zdc:eq:defect-lattice}, write each $x\in L$ uniquely as
$x=(x_0,x_1,x_2)$ with $x_i\in H_1(D_i;\mathbb Z)$.  Regard
$\mu_3$ and $\mu_4$ as their standard subgroups of $\mu_{12}$.  The inverse
on the second factor below records the candidate twist $\ell_2=-1$, rather
than a discretionary change of sign.  Define
\begin{equation}
 \begin{aligned}
 \chi_{12}:L&\longrightarrow\mu_{12},&
 \chi_{12}(x_0,x_1,x_2)&=h_1(x_1)h_2(x_2)^{-1},\\
 \chi_6:L&\longrightarrow\mu_6,&
 \chi_6(x)&=\chi_{12}(x)^2.
 \end{aligned}
 \label{zdc:eq:finite-characters}
\end{equation}

\begin{proposition}[The order-twelve and order-six reductions]
\label{zdc:prop:finite-reductions}
Put
\[
 \rho_m:\mathbb Z\longrightarrow\mu_m,
 \qquad \rho_m(k)=e^{-2\pi ik/m}.
\]
Then
\begin{equation}
 \boxed{\chi_{12}=\rho_{12}\circ\deg_F,
 \qquad \chi_6=\rho_6\circ\deg_F.}
 \label{zdc:eq:holonomy-degree}
\end{equation}
Both characters are surjective, and
\begin{equation}
 L/\ker\chi_{12}\cong\mu_{12},
 \qquad L/\ker\chi_6\cong\mu_6.
 \label{zdc:eq:finite-quotients}
\end{equation}
\end{proposition}

\begin{proof}
For
\(x=\sum_{i=0}^2(n_i a_i+r_i b_i)\),
\eqref{zdc:eq:normal-holonomies} gives, with
\(\zeta_{12}=e^{-2\pi i/12}\),
\[
 \chi_{12}(x)=\zeta_{12}^{\,4n_1-3n_2}.
\]
Equation \eqref{zdc:eq:degree-row} differs in its exponent only by
\(12n_0\), which vanishes modulo twelve.  This proves the first identity.
Squaring it is reduction of the same integer modulo six and proves the
second.  Since $\deg_F(a_1+a_2)=1$, both reductions are surjective.  The
first isomorphism theorem gives \eqref{zdc:eq:finite-quotients}.
\end{proof}

The kernels can be read without suppressing any coordinate:
\[
 \ker\chi_{12}=
 \mathbb Za_0\oplus\mathbb Zb_0\oplus
 3\mathbb Za_1\oplus\mathbb Zb_1\oplus
 4\mathbb Za_2\oplus\mathbb Zb_2,
\]
\[
 \ker\chi_6=
 \mathbb Za_0\oplus\mathbb Zb_0\oplus
 3\mathbb Za_1\oplus\mathbb Zb_1\oplus
 2\mathbb Za_2\oplus\mathbb Zb_2.
\]
The first follows by reducing $4n_1-3n_2$ modulo three and four; the
second follows by reducing it modulo three and two.  In particular, the
natural combined character has order twelve before its order-six quotient.

\subsection{Coefficient reduction is not the distinguished quotient}
\label{zdc:subsec:coefficients}

Alexander duality also gives
\[
 H^5(U;\mathbb Z)\cong\widetilde H_0(D;\mathbb Z)\cong\mathbb Z^2.
\]
The cohomological universal-coefficient sequence with coefficients
\(\mathbb Z/6\) is therefore
\[
 0\longrightarrow H^4(U;\mathbb Z)\otimes\mathbb Z/6
 \longrightarrow H^4(U;\mathbb Z/6)
 \longrightarrow
 \operatorname{Tor}_1^{\mathbb Z}(H^5(U;\mathbb Z),\mathbb Z/6)
 \longrightarrow0,
\]
whose last term is zero.  Thus the marked coefficient group is
\begin{equation}
 \boxed{H^4(U;\mathbb Z/6)\cong L/6L\cong(\mathbb Z/6)^6.}
 \label{zdc:eq:coefficient-group}
\end{equation}
If $r_6:L\to L/6L$ is coefficient reduction, then $\chi_6$ factors as
\[
 L\xrightarrow{r_6}(\mathbb Z/6)^6
 \xrightarrow{\ \overline d_F\ }\mathbb Z/6
 \xrightarrow{k\mapsto e^{-2\pi ik/6}}\mu_6,
\]
where, in the six marked residue coordinates,
\begin{equation}
 [\overline d_F]=(0,0,4,0,3,0)\pmod6.
 \label{zdc:eq:distinguished-mod-six-row}
\end{equation}
Thus
\[
 \overline d_F(\bar n_0,\bar r_0,\bar n_1,
                    \bar r_1,\bar n_2,\bar r_2)
 =4\bar n_1+3\bar n_2\in\mathbb Z/6.
\]
It is surjective because $\overline d_F(0,0,1,0,1,0)=1$ in
$\mathbb Z/6$.
Consequently
\[
 L/\ker\chi_6
 \cong (L/6L)/\ker\overline d_F
 \cong\mathbb Z/6.
\]
This is one distinguished cyclic quotient of the six-coordinate group in
\eqref{zdc:eq:coefficient-group}; it is not that group itself.

\subsection{The tensor subgroup and the genuine common-centre quotient}
\label{zdc:subsec:groups}

\begin{theorem}[Tensor subgroup and centres]
\label{zdc:thm:tensor-centres}
Tensor product defines a faithful closed homomorphism
\begin{equation}
 \iota:SU(2)\times SU(3)\hookrightarrow SU(6),
 \qquad \iota(A,B)=A\otimes B.
 \label{zdc:eq:tensor-map}
\end{equation}
Its restriction to the centres is the isomorphism
\begin{equation}
 \mu_2\times\mu_3\xrightarrow{\ \sim\ }\mu_6=Z(SU(6)),
 \qquad (\lambda,\mu)\longmapsto\lambda\mu.
 \label{zdc:eq:centre-isomorphism}
\end{equation}
The image of \(\iota\) is a proper subgroup of $SU(6)$, of real
dimension $3+8=11$.
\end{theorem}

\begin{proof}
The matrix $A\otimes B$ is unitary and
\[
 \det(A\otimes B)=(\det A)^3(\det B)^2=1,
\]
so \(\iota\) has the stated codomain.  Suppose more generally that
$A\otimes B=sI_6$.  Choose an eigenvector $w$ of $B$, with eigenvalue
$\beta$.  Then
\[
 (A\otimes B)(v\otimes w)=\beta Av\otimes w=sv\otimes w
\]
for every $v$, so $A=(s/\beta)I_2$; substitution shows that $B$ is scalar
as well.  Applying this scalar--tensor lemma with $s=1$ gives
\[
 A=\lambda I_2,\qquad B=\mu I_3,\qquad \lambda\mu=1.
\]
The determinant conditions give $\lambda^2=1$ and $\mu^3=1$.  Hence
\(\lambda=\mu^{-1}\in\mu_2\cap\mu_3=\{1\}\), proving injectivity.  A
continuous injective homomorphism from a compact group into a Hausdorff
group has closed image.

On centres the tensor matrix is $(\lambda\mu)I_6$.  Multiplication is
injective because $\mu_2\cap\mu_3=\{1\}$, and both its domain and target
have six elements, proving \eqref{zdc:eq:centre-isomorphism}.  Explicitly,
the inverse sends $s\in\mu_6$ to $(s^3,s^{-2})$.  Finally the dimension of
the image is $11<35=\dim SU(6)$, so it is proper.
\end{proof}

The isomorphisms proved above give the explicit chain
\begin{equation}
 L/\ker\chi_6\xrightarrow{\ \overline\chi_6\ }\mu_6
 \xrightarrow{\lambda\mapsto\lambda I_6}Z(SU(6)).
 \label{zdc:eq:lattice-to-centre}
\end{equation}
This is an isomorphism of finite groups.  It does not provide transition
functions for an $SU(6)$-bundle.  Also, \(\iota\) does not factor through
the quotient of $SU(2)\times SU(3)$ by its full centre: that centre maps
isomorphically, rather than trivially, to $Z(SU(6))$.

Put $\zeta_6=e^{2\pi i/6}$.  In the integer $U(1)$-charge
normalization, the common-centre generator used for the Standard-Model
global form is
\[
 (\zeta_6,-I_2,e^{2\pi i/3}I_3).
\]
If $\omega=e^{2\pi i/3}=\zeta_6^2$, then
\[
 \iota(-I_2,\omega I_3)=(-\omega)I_6=\zeta_6^{-1}I_6;
\]
the first coordinate $\zeta_6$ cancels this scalar under the
tensor--scalar map below.
This agrees with the convention in David Tong, \emph{Line Operators in the
Standard Model}, JHEP 07 (2017) 104, arXiv:1705.01853v2.  The following
matrix theorem proves the quotient directly and does not rely on a physical
interpretation of that reference.

\begin{theorem}[Common-centre quotient with $U(1)$]
\label{zdc:thm:common-centre}
Define
\begin{equation}
 \delta:\mu_6\hookrightarrow U(1)\times SU(2)\times SU(3),
 \qquad
 \delta(\xi)=(\xi,\xi^3I_2,\xi^2I_3),
 \label{zdc:eq:delta}
\end{equation}
and
\begin{equation}
 \widehat\iota:U(1)\times SU(2)\times SU(3)\longrightarrow U(6),
 \qquad
 \widehat\iota(z,A,B)=z(A\otimes B).
 \label{zdc:eq:extended-map}
\end{equation}
Then
\begin{equation}
 \ker\widehat\iota=\delta(\mu_6),
 \label{zdc:eq:extended-kernel}
\end{equation}
so the induced map is a faithful closed embedding
\begin{equation}
 \frac{U(1)\times SU(2)\times SU(3)}{\delta(\mu_6)}
 \hookrightarrow U(6).
 \label{zdc:eq:g6-embedding}
\end{equation}
Moreover,
\begin{equation}
 \boxed{\det\widehat\iota(z,A,B)=z^6,}
 \qquad
 \boxed{\operatorname{im}\widehat\iota\cap SU(6)
        =\iota(SU(2)\times SU(3)).}
 \label{zdc:eq:det-intersection}
\end{equation}
\end{theorem}

\begin{proof}
For $\xi\in\mu_6$, one has
\[
 \det(\xi^3I_2)=\xi^6=1,
 \qquad
 \det(\xi^2I_3)=\xi^6=1.
\]
Thus $\delta$ has the stated codomain.  It is central, is a homomorphism,
and is injective because its first coordinate is $\xi$.

For $\xi\in\mu_6$,
\[
 \widehat\iota(\delta(\xi))
 =\xi(\xi^3I_2\otimes\xi^2I_3)
 =\xi^6I_6=I_6,
\]
so \(\delta(\mu_6)\) lies in the kernel.  Conversely, suppose
\(z(A\otimes B)=I_6\).  The eigenvalue argument in the preceding proof,
now applied to $A\otimes B=z^{-1}I_6$, gives
\[
 A=\lambda I_2,\qquad B=\mu I_3,\qquad
 \lambda^2=1,\quad\mu^3=1,\quad z\lambda\mu=1.
\]
The Chinese-remainder isomorphism
\[
 \mu_6\longrightarrow\mu_2\times\mu_3,
 \qquad \xi\longmapsto(\xi^3,\xi^2)
\]
provides a unique $\xi\in\mu_6$ with
\(\lambda=\xi^3\) and $\mu=\xi^2$.  Then
\(z=(\lambda\mu)^{-1}=\xi\), because $\xi^{-5}=\xi$.  Hence the kernel
is exactly \(\delta(\mu_6)\).  Compactness makes the induced faithful image
closed.

The determinant formula is
\[
 \det\bigl(z(A\otimes B)\bigr)
 =z^6(\det A)^3(\det B)^2=z^6.
\]
Every tensor matrix lies in the intersection by taking $z=1$.  Conversely,
if $z(A\otimes B)$ lies in $SU(6)$, then $z^6=1$.  By
\eqref{zdc:eq:centre-isomorphism}, choose
\((\lambda,\mu)\in\mu_2\times\mu_3\) with $\lambda\mu=z$.  Then
\[
 z(A\otimes B)=(\lambda A)\otimes(\mu B),
\]
and $\lambda A\in SU(2)$, $\mu B\in SU(3)$.  This proves the reverse
inclusion in \eqref{zdc:eq:det-intersection}.
\end{proof}

Writing the quotient in \eqref{zdc:eq:g6-embedding} as $G_6$, determinant
therefore gives the exact sequence
\[
 1\longrightarrow SU(2)\times SU(3)
 \xrightarrow{\ \iota\ }G_6
 \xrightarrow{\ \det\circ\widehat\iota\ }U(1)
 \longrightarrow1.
\]
Here the middle arrow means the tensor embedding followed by the quotient
identification with $\operatorname{im}\widehat\iota$; surjectivity on the
right follows because $z\mapsto z^6$ is surjective on $U(1)$, and its kernel
is \eqref{zdc:eq:det-intersection}.

The quotient in \eqref{zdc:eq:g6-embedding} is the genuine common-centre
quotient: the $U(1)$ factor cancels the nontrivial scalar produced by the
two special-unitary centres.  The displayed tensor--scalar image in $U(6)$
is not contained in $SU(6)$; by \eqref{zdc:eq:det-intersection}, its
intersection with $SU(6)$ is exactly the tensor subgroup.  The quotient is
not isomorphic to $SU(6)$: their real dimensions are $12$ and $35$.

\subsection{Exact boundary of the result}
\label{zdc:subsec:boundary}

The construction proves the primitive fibre-degree map, its saturated
kernel, its order-twelve and order-six finite reductions, the distinction
between the full coefficient group and the selected cyclic quotient, and
the two matrix-group embeddings.  It does not produce a degree-changing
morphism from
\[
 H^4(U;\mathbb Z)\longrightarrow H^2(M;\mu_6)
\]
for a four-dimensional spacetime $M$.  In particular,
\(\chi_6\) is not, merely by its codomain, the obstruction to lifting a
common-centre quotient bundle.  No physical gauge bundle, connection,
Yang--Mills Hamiltonian, confinement statement, or mass-gap conclusion is
inferred here.
